% Figure: the law of the wall. Mean velocity u+ versus wall distance y+ in
% semilog form. The viscous sublayer follows u+ = y+; the log layer follows
% u+ = (1/kappa) ln y+ + B with kappa=0.41, B=5.0; a buffer layer bridges
% them. The x axis is logarithmic.
\begin{figure}[htbp]
\centering
\begin{tikzpicture}
\begin{semilogxaxis}[
  width=12cm, height=7cm,
  xlabel={$y^+ = y u_\tau/\nu$}, ylabel={$u^+ = u/u_\tau$},
  xmin=0.5, xmax=3000, ymin=0, ymax=25,
  legend pos=north west,
  legend cell align=left,
  grid=both, grid style={enggray!18},
  minor grid style={enggray!8},
  tick label style={font=\scriptsize},
  label style={font=\small},
  legend style={font=\scriptsize},
]
% Viscous sublayer u+ = y+, drawn 0.5..8.
\addplot[engblue, very thick, domain=0.5:8, samples=40] {x};
\addlegendentry{viscous sublayer $u^+=y^+$}
% Log law u+ = (1/0.41) ln(y+) + 5.0, drawn 30..3000.
\addplot[engred, very thick, domain=30:3000, samples=80]
  {(1/0.41)*ln(x) + 5.0};
\addlegendentry{log law $u^+=\tfrac1\kappa\ln y^+ + B$}
% DNS-like data points spanning the buffer layer and into the log region.
\addplot[engblack, only marks, mark=*, mark size=1pt] coordinates {
  (0.5,0.50) (1,1.0) (2,1.95) (4,3.80) (7,6.0) (11,8.3)
  (20,11.0) (40,13.6) (80,15.3) (150,16.9) (400,19.2)
  (1000,21.4) (2500,23.6)
};
\addlegendentry{representative data}
\end{semilogxaxis}
\end{tikzpicture}
\caption[The law of the wall]{The law of the wall for a turbulent
boundary layer, in wall units. The near-wall viscous sublayer obeys
$u^+=y^+$; beyond a buffer region the log layer obeys
$u^+=(1/\kappa)\ln y^+ + B$ with $\kappa\approx0.41$ and $B\approx5.0$.
This two-region structure is one of the few quantitative universals of
wall-bounded turbulence and underlies near-wall meshing rules in CFD.}
\label{fig:vol03ch12:law-of-wall}
\end{figure}
